% ****************************************************************************************************
% Capítulo 7: Antipsicóticos na Prática Clínica
% ****************************************************************************************************
\chapter{Antipsicóticos na Prática Clínica}
\label{ch:antipsicoticos}

\begin{flushright}
\textit{Ferramenta Terapêutica ou Armadilha Medicalizante?}
\end{flushright}

\vspace{1em}

% ****************************************************************************************************
\section{A Dupla Face dos Antipsicóticos}
\label{sec:dupla-face}
% ****************************************************************************************************

Os antipsicóticos ocupam posição singular no arsenal terapêutico da psiquiatria contemporânea. Como observa Rose (2019), estes fármacos representam simultaneamente conquista técnica inegável e vetor potencial de medicalização indiscriminada. Desde a introdução da clorpromazina em 1952 --- evento que Healy (2002) denominou \enquote{revolução psicofarmacológica} ---, a psiquiatria oscila entre a celebração de sua capacidade de aliviar sofrimento psicótico genuíno e a crítica de seu uso como instrumento de controle social.

Esta tensão reflete o que Foucault (1961/2006) identificou como ambiguidade constitutiva da psiquiatria moderna: disciplina que, ao mesmo tempo em que promete libertação do sofrimento mental, opera como dispositivo de poder que delimita fronteiras entre o normal e o patológico.

O campo contemporâneo enfrenta o que Hacking (1999) denominou \enquote{nominalismo dinâmico} aplicado à psicofarmacologia: as categorias diagnósticas que justificam prescrições não apenas descrevem estados mentais preexistentes, mas ativamente moldam a experiência subjetiva que pretendem tratar. Como alertam Moncrieff e Cohen (2006) em sua crítica ao modelo centrado na doença, a suposição de que antipsicóticos corrigem desequilíbrios neuroquímicos específicos permanece empiricamente contestável.

Este capítulo propõe uma análise epistemologicamente rigorosa do uso de antipsicóticos, diferenciando indicações genuinamente terapêuticas de aplicações que representam o que Conrad (2007) caracteriza como \enquote{medicalização}.

% ****************************************************************************************************
\section{O que Define Indicação Genuína?}
\label{sec:indicacao-genuina}
% ****************************************************************************************************

\subsection{A Fenomenologia do Sintoma Psicótico Verdadeiro}

A questão da indicação adequada remete ao problema mais fundamental da definição de psicose genuína. Jaspers (1913/1997) estabeleceu distinção seminal entre \enquote{desenvolvimento da personalidade} --- transformações compreensíveis da experiência --- e \enquote{processo} --- alterações que introduzem descontinuidade radical na biografia do sujeito.

Como observa Stanghellini (2016), o sintoma psicótico genuíno caracteriza-se pelo que denomina \enquote{perturbação da ipseidade} --- alteração na estrutura básica da experiência de si mesmo que precede e fundamenta manifestações sintomáticas específicas.

Os antipsicóticos encontram sua indicação genuína precisamente nestas condições caracterizadas por:

\begin{itemize}
\item \textbf{Alterações estruturais da experiência}: Não meramente conteúdos de pensamento inusuais, mas transformações no que Ratcliffe (2008) denomina \enquote{sentimentos existenciais} --- a estrutura básica através da qual experimentamos o mundo como real.

\item \textbf{Impacto funcional pervasivo e persistente}: Comprometimento significativo da capacidade de \enquote{integração narrativa}, como descrevem Bolton e Hill (2004).

\item \textbf{Base neurobiológica identificável}: Alterações consistentes em conectividade frontotemporal, como documenta Meyer-Lindenberg (2010).
\end{itemize}

\subsection{A Esquizofrenia como Paradigma}

A esquizofrenia representa o paradigma de condição para a qual antipsicóticos demonstram eficácia robusta. Meta-análises documentam que antipsicóticos reduzem significativamente sintomas positivos e reduzem em aproximadamente 60\% o risco relativo de recaída.

Entretanto, como observa Moncrieff (2013), estes dados devem ser interpretados com cautela. A demonstração de eficácia não implica necessariamente que antipsicóticos \enquote{corrijam} um desequilíbrio neuroquímico subjacente. Pode refletir o que ela denomina \enquote{modelo centrado no fármaco} --- os antipsicóticos produzem estados mentais alterados que, incidentalmente, atenuam manifestações psicóticas.

\subsection{Quadros Transitórios e o Papel do Aripiprazol}

Nem todo sintoma psicótico justifica tratamento prolongado. Como observa van Os (2009), experiências semelhantes a alucinações e ideação paranoide ocorrem em proporção significativa da população geral sem configurar patologia.

Para quadros caracterizados por descontinuidades mais sutis --- o que denominamos \enquote{disrupturas cognitivas} com colorido psicótico ---, o aripiprazol oferece perfil farmacológico diferenciado. Como agonista parcial dopaminérgico D2, demonstra efeito \enquote{estabilizador} --- modulando a neurotransmissão sem o bloqueio massivo característico de antipsicóticos convencionais.

% ****************************************************************************************************
\section{O Problema do Uso Inadequado}
\label{sec:uso-inadequado}
% ****************************************************************************************************

\subsection{A Medicalização do Sofrimento Existencial}

A expansão contínua das indicações de antipsicóticos representa manifestação paradigmática do que Illich (1976/2010) identificou como \enquote{expropriação da saúde} --- processo pelo qual dimensões da experiência humana são progressivamente colonizadas pela racionalidade biomédica.

Esta tendência manifesta-se particularmente em duas populações vulneráveis: pessoas com Transtorno do Espectro Autista (TEA) e pessoas com Deficiência Intelectual (DI). O programa STOMP do NHS Inglaterra documenta que 17,1\% dos adultos com DI recebem prescrições de antipsicóticos, frequentemente para \enquote{comportamentos desafiadores} --- indicação com base evidencial consideravelmente mais frágil.

\subsection{Diagnóstico Errôneo de Esquizofrenia}

A atribuição precipitada de diagnóstico de esquizofrenia representa uma das formas mais consequentes de uso inadequado. Como observa Frances (2013), a psiquiatria contemporânea sofre de \enquote{inflação diagnóstica}.

Condições frequentemente confundidas com esquizofrenia incluem:

\begin{itemize}
\item \textbf{Estados dissociativos}: O que Janet (1889/1973) descreveu como \enquote{desagregação psíquica}.

\item \textbf{Transtornos de personalidade}: Particularmente o transtorno borderline, com seus \enquote{episódios micropsicóticos} transitórios.

\item \textbf{Intoxicação e abstinência de substâncias}: Estados que podem mimetizar esquizofrenia mas remitem com abstinência.

\item \textbf{Condições médicas gerais}: Encefalopatias, endocrinopatias e epilepsias do lobo temporal.
\end{itemize}

\subsection{Autismo e a Confusão com Psicose}

Pessoas autistas frequentemente apresentam estrutura de linguagem que pode ser erroneamente interpretada como \enquote{desorganização do pensamento}. Como argumenta Milton (2012) em sua teoria do \enquote{problema da dupla empatia}, dificuldades comunicativas representam não déficit unilateral, mas incompatibilidade bidirecional entre estilos cognitivos distintos.

Medicalizar esta diferença com antipsicóticos constitui o que Goffman (1963/1986) caracterizou como \enquote{estigmatização secundária}.

% ****************************************************************************************************
\section{Antipsicóticos no Espectro Terapêutico}
\label{sec:espectro-terapeutico}
% ****************************************************************************************************

\subsection{O Princípio da Graduação Interventiva}

A prática psiquiátrica eticamente orientada opera segundo o que podemos denominar \enquote{princípio da graduação interventiva} --- a preferência por intervenções menos invasivas quando igualmente eficazes. O imperativo de não-maleficência demanda que intervenções com potencial iatrogênico significativo sejam reservadas para situações em que alternativas menos arriscadas demonstraram-se insuficientes.

\subsection{Intervenções Psicossociais: Evidências}

A meta-análise de rede de Salahuddin et al. (2024) demonstrou que Terapia Cognitivo-Comportamental para psicose (TCCp) foi significativamente superior ao tratamento padrão na redução de sintomas globais, com tamanho de efeito comparável ao de antipsicóticos em várias análises.

A TCCp opera através de mecanismos que incluem: normalização de experiências psicóticas; redução de processos que mantêm sintomas; desenvolvimento de explicações alternativas para experiências anômalas; e fortalecimento de estratégias de coping.

\subsection{Integração de Modalidades}

A prática baseada em evidências não opõe intervenções farmacológicas a psicossociais, mas busca sua integração ótima. Como demonstra Bighelli et al. (2021), intervenções familiares, psicoeducação e TCCp reduzem significativamente taxas de recaída quando combinadas a farmacoterapia --- efeito sinérgico que nenhuma modalidade isolada alcança.

% ****************************************************************************************************
\section{Reflexões Éticas e Práticas}
\label{sec:reflexoes-eticas}
% ****************************************************************************************************

\subsection{O Imperativo da Decisão Compartilhada}

A prescrição de antipsicóticos envolve decisões que afetam profundamente a experiência subjetiva, o funcionamento cognitivo e a qualidade de vida. Como argumenta Deegan (2007), pacientes não são recipientes passivos de intervenções técnicas, mas agentes ativos na construção de suas trajetórias de recuperação.

O modelo de \enquote{decisão compartilhada} propõe que médico e paciente sejam reconhecidos como experts em domínios complementares --- o clínico em evidências científicas, o paciente em seus valores e experiência vivida.

\subsection{Preservação da Subjetividade}

Estudos qualitativos documentam consistentemente que efeitos adversos subjetivos --- sedação, embotamento afetivo, anedonia --- constituem os principais determinantes de não-adesão ao tratamento.

O objetivo do tratamento não deve ser a \enquote{normalização} do paciente segundo padrões externamente definidos, mas o que Fulford (1989) denomina \enquote{restauração da capacidade de ação intencional}.

\subsection{Educação e Consentimento Informado}

O consentimento verdadeiramente informado para tratamento antipsicótico requer que o paciente compreenda não apenas benefícios potenciais, mas também:

\begin{itemize}
\item \textbf{Efeitos adversos}: Riscos metabólicos, neurológicos e subjetivos.
\item \textbf{Alternativas}: Intervenções psicossociais com base evidencial comparável para determinados desfechos.
\item \textbf{Incertezas}: O fato de que mecanismos de ação permanecem incompletamente compreendidos.
\end{itemize}

% ****************************************************************************************************
\section{Propostas para Uso Racional}
\label{sec:uso-racional}
% ****************************************************************************************************

\subsection{Avaliação Fenomenológica Criteriosa}

Antes de prescrever antipsicóticos, o clínico deve realizar o que Minkowski (1927/1970) denominou \enquote{diagnóstico por penetração} --- compreensão empática da estrutura alterada da experiência, não meramente catalogação de sintomas.

\subsection{Monitoramento Sistemático}

O monitoramento deve incluir:

\begin{itemize}
\item \textbf{Parâmetros metabólicos}: Peso, glicemia, perfil lipídico, pressão arterial.
\item \textbf{Função cognitiva}: Avaliação neuropsicológica periódica.
\item \textbf{Experiência subjetiva}: Escalas de qualidade de vida e entrevistas qualitativas.
\end{itemize}

\subsection{Revisão Periódica e Possibilidade de Descontinuação}

A prescrição não deve ser presumidamente indefinida. Como propõem Horowitz e Taylor (2024), a descontinuação deve seguir princípio de \enquote{redução hiperbólica} --- reduções proporcionalmente menores à medida que se aproxima de doses mais baixas.

% ****************************************************************************************************
\section{Entre a Ferramenta e a Armadilha}
\label{sec:ferramenta-armadilha}
% ****************************************************************************************************

Os antipsicóticos representam ferramenta terapêutica indispensável para condições psicóticas genuínas. A esquizofrenia, transtornos delirantes e psicoses afetivas frequentemente requerem intervenção farmacológica para restabelecer o que Jaspers (1913/1997) denominou \enquote{contato com a realidade}.

Entretanto, o uso criterioso requer diferenciação rigorosa entre indicações genuínas e aplicações que representam medicalização inadequada. A prática clínica eticamente orientada demanda: avaliação fenomenológica profunda; integração de intervenções farmacológicas com abordagens psicossociais; decisão compartilhada que respeite autonomia; monitoramento sistemático; e revisão periódica da necessidade de continuação.

Como observa Kleinman (1988), o sofrimento mental não se reduz a disfunção neuroquímica corrigível farmacologicamente, mas envolve dimensões experienciais, relacionais e significativas que requerem abordagem compreensiva. Os antipsicóticos, quando prescritos criteriosamente, podem facilitar esta abordagem; quando prescritos indiscriminadamente, podem obstaculizá-la, substituindo compreensão por supressão e transformação por normalização.

% ****************************************************************************************************
% Referências do Capítulo
% ****************************************************************************************************
\vspace{2em}
\begin{small}
\noindent\textsc{Referências}

\vspace{1em}
\noindent Bighelli, I., et al. (2021). Psychosocial interventions for relapse prevention in schizophrenia. \textit{Lancet Psychiatry}, 8(11), 969--980.

\noindent Conrad, P. (2007). \textit{The medicalization of society}. Johns Hopkins University Press.

\noindent Deegan, P. E. (2007). The lived experience of using psychiatric medication. \textit{Psychiatric Rehabilitation Journal}, 31(1), 62--69.

\noindent Frances, A. (2013). \textit{Saving normal}. William Morrow.

\noindent Healy, D. (2002). \textit{The creation of psychopharmacology}. Harvard University Press.

\noindent Jaspers, K. (1997). \textit{General psychopathology}. (Orig. 1913). Johns Hopkins University Press.

\noindent Kleinman, A. (1988). \textit{The illness narratives}. Basic Books.

\noindent Moncrieff, J. (2013). \textit{The bitterest pills}. Palgrave Macmillan.

\noindent Rose, N. (2019). \textit{Our psychiatric future}. Polity.

\noindent Salahuddin, N. H., et al. (2024). Psychological interventions for treatment-resistant schizophrenia. \textit{Lancet Psychiatry}, 11(7), 545--553.
\end{small}
